\ssr{Введение}

Цели данной практической работы:
\begin{enumerate}
	\item изучить этапы загрузки операционных систем семейства GNU/Linux;
	\item освоить методы анализа процесса загрузки;
	\item получить навыки настройки параметров загрузчика, диагностики проблем и оптимизации времени старта системы.
\end{enumerate}

\ssr{Выполнение}

\ssection{Задание 1. Карта загрузки системы}

Цель: понять последовательность этапов загрузки Linux, определить основные компоненты и проанализировать параметры, переданные ядру.

Задачи:
\begin{enumerate}
	\item определить систему инициализации;
	\item изучить параметры загрузки ядра и объяснить их назначение;
	\item оценить время загрузки;
	\item посмотреть загруженные модули ядра.
\end{enumerate}

\begin{figure}[!htbp]
	\centering
	\includegraphics[width=0.6\linewidth]{"Screenshot From 2026-10-05 13-50-32"}
	\caption{Вывод системы инициализации и загруженных модулей ядра}
\end{figure}
\FloatBarrier

\begin{figure}[!htbp]
	\centering
	\includegraphics[width=0.6\linewidth]{"Screenshot 2026-10-07 133738"}
	\caption{Фрагмент boot.svg}
\end{figure}
\FloatBarrier
\clearpage
Вопросы для отчета:
\begin{enumerate}
	\item \textit{Перечислите этапы загрузки в правильном порядке (от включения питания до появления приглашения входа).}\newline 
	BIOS/UEFI $\rightarrow$ GRUB $\rightarrow$ ядро $\rightarrow$ initramfs $\rightarrow$ systemd $\rightarrow$ target $\rightarrow$ приглашение входа.
	\item \textit{Какой процесс запускается первым после ядра? Какова его роль?}\newline 
	Процесс systemd (PID 1). 
	Его роль -- инициализация системы, запуск юнитов и сервисов, управление целями (targets), монтирование ФС и т.д.
	\item \textit{Сколько времени заняла загрузка вашей системы (kernel + userspace)?}\newline 
	11.921 секунд.
	\item \textit{Запишите содержимое /proc/cmdline. Что означают параметры quiet, splash, ro? Какой UUID у вашей корневой файловой системы? Зачем используется UUID вместо имени устройства (например, /dev/sda1)?}\newline 
	Параметры: quiet -- подавление подробных сообщений ядра; splash -- показ графической заставки; ro -- монтирование корневой ФС в режиме только для чтения на этапе initramfs.
	UUID моей системы -- ?.
	UUID используется благодаря его стабильности, так как имена могут меняться при подключении дисков.
	\item \textit{Какие модули ядра загружены первыми? (из вывода lsmod)}\newline
	Первыми в выводе lsmod отображаются модули, загруженные последними (обратный порядок).
\end{enumerate}

\ssection{Задание 2. Исследование загрузчика GRUB}

Цель: изучить конфигурацию загрузчика GRUB, научиться находить доступные ядра и параметры загрузки.

Задачи:
\begin{enumerate}
	\item понять структуру конфигурационного файла GRUB;
	\item определить все доступные пункты меню (включая подменю);
	\item определить настройки по умолчанию.
\end{enumerate}
\clearpage

\begin{figure}[!htbp]
	\centering
	\includegraphics[width=0.6\linewidth]{"Screenshot From 2026-10-05 13-54-28"}
	\includegraphics[width=0.6\linewidth]{"Screenshot From 2026-10-05 13-54-57"}	
	\caption{Фрагмент файла конфигурации GRUB}
\end{figure}
\FloatBarrier

\begin{figure}[!htbp]
	\centering
	\includegraphics[width=0.6\linewidth]{"Screenshot From 2026-10-05 13-55-20"}
	\includegraphics[width=0.6\linewidth]{"Screenshot From 2026-10-05 13-56-43"}
	\caption{Фрагмент файла конфигурации GRUB, отображающий menuentry}
\end{figure}
\FloatBarrier
\clearpage

\begin{figure}[!htbp]
	\centering
	\includegraphics[width=0.6\linewidth]{"Screenshot From 2026-10-05 13-58-08"}
	\includegraphics[width=0.6\linewidth]{"Screenshot From 2026-10-05 13-59-52"}
	\caption{Параметры конфигурации GRUB}
\end{figure}
\FloatBarrier

Вопросы для отчета:
\begin{enumerate}
	\item \textit{Сколько различных пунктов меню (включая подменю) вы обнаружили в файле grub.cfg?}\newline 
	12 пунктов меню.
	\item \textit{Сколько ядер доступно для загрузки в вашей системе? Укажите их версии.}\newline 
	Одно ядро: 7.0.0-31-generic (файлы /boot/vmlinuz, /boot/vmlinuz.old, /boot/vmlinuz-7.0.0-31-generic -- это одно и то же ядро).
	\item \textit{Какое ядро загружается по умолчанию? Как это определено в /etc/default/grub?}\newline 
	Пункт по умолчанию -- первый (Ubuntu с ядром 7.0.0-31-generic), так как GRUB\_DEFAULT=0.
	Параметром GRUB\_DEFAULT: 0 -- первый пункт меню, saved -- последнее загруженное, либо точное имя пункта.
	\item \textit{Какой параметр в файле /etc/default/grub отвечает за время отображения меню? Какое значение установлено у вас?}\newline 
	GRUB\_TIMEOUT (в секундах). 
	% Также используется GRUB\_TIMEOUT\_STYLE (menu/countdown/hidden).
	\item \textit{Что произойдет, если установить GRUB\_TIMEOUT=0?}\newline  
	Меню не отображается, сразу грузится пункт по умолчанию. Для вызова меню нужно удерживать Shift (BIOS) или Esc (UEFI).
\end{enumerate}

\ssection{Задание 3. Анализ времени загрузки}

Цель: научиться определять, какие службы замедляют старт системы.

Задачи:
\begin{enumerate}
	\item получить общее время загрузки;
	\item найти самые «медленные» сервисы;
	\item построить цепочку критических зависимостей.
\end{enumerate}

\begin{figure}[!htbp]
	\centering
	\includegraphics[width=0.6\linewidth]{"Screenshot From 2026-10-05 15-56-40.png"}
	\includegraphics[width=0.6\linewidth]{"Screenshot From 2026-10-05 15-57-35.png"}
	\caption{Вывод времени загрузки и цепочки критических зависимостей}
\end{figure}
\FloatBarrier

Вопросы для отчета:
\begin{enumerate}
	\item \textit{Какой сервис грузится дольше всех? Сколько времени это занимает?}\newline 
	Сервис dev-rfkill.device, 6.868 секунд. 
	\item \textit{Что означает состояние active (running) у сервиса?}\newline 
	Что сервис запущен и работает в данный момент.
	\item \textit{Какой target отвечает за графический интерфейс? Как проверить, активен ли он?}\newline 
	graphical.target.
	Проверка: systemctl is-active graphical.target.
	\item \textit{Как переключиться в текстовый режим без перезагрузки?}\newline 
	Командой sudo systemctl isolate multi-user.target (графические приложения закроются).
\end{enumerate}

\ssection{Задание 4. Исследование Initramfs}

Цель: понять назначение временной корневой файловой системы (initramfs).

Задачи:
\begin{enumerate}
	\item найти файл initramfs для текущего ядра;
	\item оценить его размер и дату создания;
	\item изучить содержимое initramfs.
\end{enumerate}

\begin{figure}[!htbp]
	\centering
	\includegraphics[width=0.6\linewidth]{"Screenshot From 2026-10-05 15-58-58.png"}
	\caption{Вывод содержимого файла initramfs}
\end{figure}
\FloatBarrier

Вопросы для отчета:
\begin{enumerate}
	\item \textit{Какой размер у вашего initramfs? (в мегабайтах)}\newline 
	42 Мб (команда ls -lh /boot/initrd.img-7.0.0-31-generic).
	\item \textit{Для чего нужен initramfs?}\newline 
	Временная ФС с драйверами и утилитами для монтирования реальной корневой ФС.
	\item \textit{Какие драйверы (модули) присутствуют в initramfs? Приведите 2–3 примера.}\newline 
	По выводу lsinitramfs модулей .ko не найдено (grep -E ".ko\$" пуст); в initramfs присутствуют микрокод kernel/x86/microcode/GenuineIntel.bin, скрипты и конфигурация (etc/cmdline.d/20-root-dev.conf).
	\item \textit{Что произойдёт, если удалить initramfs?}\newline 
	Система не загрузится, если драйверы корня не встроены в ядро.
\end{enumerate}

\ssection{Задание 5. Мониторинг процесса загрузки (логи)}

Цель: научиться читать и анализировать сообщения, генерируемые при загрузке.

Задачи:
\begin{enumerate}
	\item просмотреть логи текущей загрузки;
	\item найти сообщения ядра;
	\item обнаружить возможные ошибки.
\end{enumerate}

\begin{figure}[!htbp]
	\centering
	\includegraphics[width=0.6\linewidth]{"Screenshot From 2026-10-05 15-59-46.png"}
	\includegraphics[width=0.6\linewidth]{"Screenshot From 2026-10-19 12-26-11.png"}
	\includegraphics[width=0.6\linewidth]{"Screenshot From 2026-10-19 12-26-46.png"}
	\includegraphics[width=0.6\linewidth]{"Screenshot From 2026-10-05 16-00-23.png"}
	\caption{Вывод логов текущей загрузки}
\end{figure}
\FloatBarrier

Вопросы для отчета:
\begin{enumerate}
	\item \textit{В какое время началась текущая загрузка системы?}\newline 
	2026-10-05 в 14:59:41.
	\item \textit{Были ли ошибки при загрузке? Если да, то какие?}\newline 
	Да: 
	piix4\_smbus — SMBus Host Controller not enabled!; 
	gkr-pam: couldn't unlock the login keyring; /dev/sr0;
	/dev/sr1: Can't open blockdev.
	\item \textit{Какой сервис, судя по логам, запустился первым? (после ядра)}\newline 
	systemd.
\end{enumerate}

\ssection{Задание 6. Сравнение состояний системы}

Цель: проанализировать состояние уже загруженной системы.

Задачи:
\begin{enumerate}
	\item подсчитать загруженные модули ядра;
	\item посмотреть смонтированные файловые системы;
	\item изучить сетевые интерфейсы и запущенные процессы.
\end{enumerate}

\begin{figure}[!htbp]
	\centering
	\includegraphics[width=0.6\linewidth]{"Screenshot From 2026-10-05 16-00-57.png"}
	\caption{Вывод количества загруженных модулей ядра и списка смонтированных файловых систем}
\end{figure}
\FloatBarrier

Вопросы для отчета:
\begin{enumerate}
	\item \textit{Сколько модулей ядра загружено в системе?}\newline 
	64 модуля (lsmod | wc -l дал 65 строк, из них 1 -- заголовок).
	\item \textit{Какая файловая система используется для корневого каталога (`/`)?}\newline 
	ext4 (/dev/sda2 на /).
	\item \textit{Какие сетевые интерфейсы активны? Есть ли у них IP-адреса?}\newline
	lo (127.0.0.1, ::1) и ens33 (192.168.237.130/24, fe80::20c:29ff:fe4c:ddf1/64). Оба активны и имеют IP-адреса.
	\item \textit{Сколько процессов работает в системе? (примерно)}\newline 
	Около 391 процесса.
\end{enumerate}

\ssection{Задание 7. Эксперимент с уровнями выполнения (targets)}

Цель: понять разницу между различными состояниями системы в systemd.

Задачи:
\begin{enumerate}
	\item изучить доступные targets;
	\item узнать текущий target по умолчанию;
	\item сравнить multi-user.target и graphical.target.
\end{enumerate}

\begin{figure}[!htbp]
	\centering
	\includegraphics[width=0.6\linewidth]{"Screenshot From 2026-10-05 16-01-32.png"}
	\caption{Вывод юнитов targets}
\end{figure}
\FloatBarrier

\clearpage
\begin{figure}[!htbp]
	\centering
	\includegraphics[width=0.6\linewidth]{"Screenshot From 2026-10-05 16-01-47.png"}
	\caption{Вывод зависимостей для multi-user.target и graphical.target}
\end{figure}
\FloatBarrier

\begin{figure}[!htbp]
	\centering
	\includegraphics[width=0.6\linewidth]{"Screenshot From 2026-10-05 16-02-07.png"}
	\caption{Вывод доступных targets}
\end{figure}
\FloatBarrier

Вопросы для отчета:
\begin{enumerate}
	\item \textit{В каком target работает ваша система сейчас?}\newline 
	Определяется командой systemctl get-default (по умолчанию) или systemctl list-units --type=target --state=active
	\item \textit{Какие сервисы есть в graphical.target, но отсутствуют в multi-user.target?}\newline 
	Сравнение выводов systemctl list-dependencies graphical.target и multi-user.target; обычно display-manager.service (gdm и т.п.).
	\item \textit{Как переключиться в текстовый режим (multi-user) без перезагрузки?}\newline 
	Выполнить команду (с предупреждением о закрытии графических приложений): sudo systemctl isolate multi-user.target
\end{enumerate}

%\ssection{Дополнительное задание}

\ssr{Вывод}

В ходе практической работы были:
\begin{itemize}
	\item изучены этапы загрузки операционных систем семейства GNU/Linux;
	\item освоены методы анализа процесса загрузки;
	\item получены навыки настройки параметров загрузчика, диагностики проблем и оптимизации времени старта системы.
\end{itemize}

\ssr{Контрольные вопросы}

\begin{enumerate}
	\item \textit{Назовите основные этапы загрузки Linux.}\newline 
	BIOS/UEFI $\rightarrow$ GRUB $\rightarrow$ ядро $\rightarrow$ initramfs $\rightarrow$ systemd $\rightarrow$ target $\rightarrow$ приглашение входа.
	\item \textit{В чём различие между BIOS и UEFI?}\newline 
	BIOS использует MBR (первый сектор диска), UEFI использует GPT и раздел ESP с файлами .efi. UEFI поддерживает Secure Boot, большие диски, быстрее загружается.
	\item \textit{Какова роль GRUB в процессе загрузки?}\newline 
	Загрузка ядра и initramfs, передача параметров ядру, предоставление меню выбора ОС/ядра.
	\item \textit{Что такое initramfs и для чего он нужен?}\newline 
	Временная ФС с драйверами и утилитами для монтирования реального корня. Нужна для загрузки драйверов, не встроенных в ядро.
	\item \textit{Какая информация содержится в файле /proc/cmdline?}\newline 
	Параметры, переданные ядру загрузчиком.
	\item \textit{Как узнать, какая система инициализации используется?}\newline 
	Команда ps -p 1 -o comm=.
	\item \textit{Перечислите основные команды systemd-analyze.}\newline 
	time (общее время), blame (время по сервисам), critical-chain (критическая цепочка), plot > boot.svg (график).
	\item \textit{Как посмотреть логи только текущей загрузки?}\newline 
	Команда journalctl -b.
	\item \textit{Что такое target в systemd? Приведите примеры.}\newline 
	Целевое состояние системы. Примеры: multi-user.target (без графики), graphical.target (с графикой), rescue.target (восстановление), reboot.target, poweroff.target.
	\item \textit{Как добавить параметр ядра временно (при загрузке) и постоянно?}\newline 
	Временно: в меню GRUB нажать e, добавить параметр в строку linux, загрузиться. Постоянно: добавить в GRUB\_CMDLINE\_LINUX\_DEFAULT в /etc/default/grub, затем sudo update-grub.
\end{enumerate}